\begin{definition}[The sub-elliptic setting]
A \emph{Carnot setting} is a Carnot norm datum (\noderef{CarnotNorm}) together with a map
$\sigma:\mathbb R^n\to\mathbb R^{n\times m}$, $\sigma(x)=[\sigma^1(x),\ldots,\sigma^m(x)]$
--- the horizontal frame, written $X_i=\sigma^i(x)\cdot\nabla$ --- subject to the following
conditions. Here $\mathcal Xu(x)=\mathrm{hgrad}\ \sigma\ x\ \nabla u(x)$ is the horizontal
gradient of \noderef{hgrad}, $\mathcal X\mathcal X^{s}$ is the symmetrized horizontal
second derivative of \noderef{hhessS}, a \emph{second-order jet} is as in
\noderef{TwiceDiffAt}, and $L_h(y)=h*y$ denotes left translation.
\begin{enumerate}
\item Each $\sigma^i:\mathbb R^n\to\mathbb R^n$ is $C^\infty$.
\item For every $x\in\mathbb R^n$ the vectors $\sigma^1(x),\ldots,\sigma^m(x)$ are linearly
  independent.
\item (H\"ormander's rank condition (1.6).) For every $x\in\mathbb R^n$,
  \[
   \mathrm{span}\,\big\{\,V(x)\ :\ V\in\mathrm{LieHull}\ \sigma\,\big\}=\mathbb R^n,
  \]
  with $\mathrm{LieHull}\ \sigma$ the Lie hull of the frame (\noderef{LieHull}).
\item (The Taylor formula (2.18).) If $u:\mathbb R^n\to\mathbb R$ is differentiable at
  $x$ with differential $D$, then $u$ is Pansu differentiable at $x$
  (\noderef{HasPansuGradAt}) with horizontal differential
  $\mathrm{hgrad}\ \sigma\ x\ D$; that is,
  $u(x*h)=u(x)+\langle\mathcal Xu(x),\pi_1(h)\rangle+o(\|h\|)$ as $h\to0$.
\item (Left-invariance of the horizontal derivatives.) For all $h,x\in\mathbb R^n$ and
  every $\varphi:\mathbb R^n\to\mathbb R$ with a second-order jet $(D,B)$ at $h*x$, the
  function $\varphi\circ L_h$ has a second-order jet $(D',B')$ at $x$ satisfying
  \[
   \mathrm{hgrad}\ \sigma\ x\ D'=\mathrm{hgrad}\ \sigma\ (h*x)\ D,
   \qquad
   \mathrm{hhessS}\ \sigma\ x\ D'\ B'=\mathrm{hhessS}\ \sigma\ (h*x)\ D\ B .
  \]
\item (The form of the $X_i$ on functions of the first layer.) For every
  $\psi:\mathbb R^m\to\mathbb R$ with a second-order jet $(E,C)$ at $\pi_1(x)$, the
  function $\psi\circ\pi_1$ has a second-order jet $(D,B)$ at $x$ satisfying
  \[
   \mathrm{hgrad}\ \sigma\ x\ D=\big(E(e_i)\big)_{i=1}^{m},
   \qquad
   \mathrm{hhessS}\ \sigma\ x\ D\ B=\big(C(e_i,e_j)\big)_{i,j=1}^{m},
  \]
  with $e_1,\ldots,e_m$ the standard basis of $\mathbb R^m$.
\end{enumerate}

Items (1)--(3) are the hypotheses on the vector fields in Section Section~1.2: in
exponential coordinates $X_i=\partial_{x_{i,1}}+\sum_{j\ge2}\sum_kb^i_{k,j}(x)\partial_{x_{k,j}}$
with polynomial $b^i_{k,j}$, so the $\sigma^i$ are smooth and, their first-layer components
forming the identity matrix, pointwise linearly independent --- in particular
$\sigma(x)\neq0$ --- and $X_1,\ldots,X_m$, being a basis of the first layer of a
stratification, satisfy (1.6). Item (4) is (2.18), which the paper derives from the
Euclidean expansion of $u$ at $x$ together with
$x*h-x=\sum_i\pi_1(h)_iX_i(x)+O(\|h\|^2)$.

Items (5) and (6) are the two structural facts stated at the level of jets, which is where
the development uses them. Item (5) is left-invariance of the $X_i$: since
$X_i(\varphi\circ L_h)=(X_i\varphi)\circ L_h$ and
$X_jX_i(\varphi\circ L_h)=(X_jX_i\varphi)\circ L_h$, and since $L_h$ is a polynomial
diffeomorphism, a second-order jet $(D,B)$ of $\varphi$ at $h*x$ transports to the jet
$D'=D\circ DL_h(x)$, $B'(u,v)=B(DL_h(x)u,DL_h(x)v)+D(D^2L_h(x)(u,v))$ of
$\varphi\circ L_h$ at $x$, whose horizontal gradient and symmetrized horizontal second
derivative at $x$ are those of $\varphi$ at $h*x$. This is what makes Lemma Lemma~2.9
--- left translates of sub-solutions are sub-solutions --- a bookkeeping statement. Item (6)
is the displayed computation in the proof of Theorem Theorem~1.3(iii): because
$X_i=\partial_{x_{i,1}}+\sum_{j\ge2}\sum_kb^i_{k,j}\partial_{x_{k,j}}$ kills the dependence
on the higher layers, $X_i(\psi\circ\pi_1)=(\partial_i\psi)\circ\pi_1$, and iterating,
$X_jX_i(\psi\circ\pi_1)=(\partial_j\partial_i\psi)\circ\pi_1$; hence
$\mathcal X(\psi\circ\pi_1)=\nabla\psi\circ\pi_1$ and
$\mathcal X\mathcal X^{s}(\psi\circ\pi_1)=D^2\psi\circ\pi_1$, which is the displayed
identity $\mathcal L(\psi\circ\pi_1)=(\mathcal L^{\mathrm{eucl}}\psi)\circ\pi_1$ once the
operator is applied. Both items are stated as the existence of a transported jet with the
prescribed horizontal data, which is all the proofs use; the transported jet is in fact
unique, as second-order jets always are.

The formalization note (item 3 of the note after Theorem Theorem~1.3) permits exactly these
facts of Section Section~2.4 to be built into the definition of the setting, provided
every Carnot group in those coordinates satisfies the definition, and each condition above
is one of them. Deliberately absent is any statement about sub- or super-solutions of
$\mathcal L$: the strong maximum principle, and a fortiori any comparison or strong
comparison principle, is a cited result carried as a separate hypothesis, not a property of
the setting. For the Euclidean additive group on $\mathbb R^m$ ($n=m$, $r=1$, $x*y=x+y$,
$\|x\|=|x|$, $\pi_1=\mathrm{id}$, $\exp_{\mathrm h}=\mathrm{id}$) the frame is the constant
one $\sigma^i(x)=e_i$: then (1) and (2) are clear, the Lie hull consists of the $e_i$ and
the zero field so (3) holds because $e_1,\ldots,e_m$ span $\mathbb R^m$,
$\mathcal M(x,D)=0$ so $\mathcal Xu=\nabla u$ and
$\mathcal X\mathcal X^{s}u=D^2u$, (4) is the definition of the Fr\'echet derivative,
and (5), (6) hold with $(D',B')=(D,B)$ and $(D,B)=(E,C)$ since $DL_h=I$, $D^2L_h=0$ and
$\pi_1=\mathrm{id}$.
\end{definition}
